\documentclass[11pt]{article}
\usepackage[utf8]{inputenc}
\usepackage[T1]{fontenc}
\usepackage{lmodern}
\usepackage[margin=1in]{geometry}
\usepackage{csquotes}
\usepackage[style=authoryear,backend=biber,sorting=nyt,maxbibnames=99,giveninits=true,uniquename=init]{biblatex}
\usepackage[hidelinks]{hyperref}
\addbibresource{references.bib}
\begin{document}
\section*{Model Performance Evaluation and Comparison}

\subsection{Overview}

Evaluating and comparing the accuracy of predictive tools is a recurring
requirement across the tile-drainage literature, spanning physically based
drainage simulators, analytical water-table equations, statistical and
machine-learning surrogates, geostatistical interpolators, remote-sensing
classifiers, and field measurement or sampling protocols. The evidence base
assembled here converges on several well-supported conclusions. First, the
choice of goodness-of-fit metric or objective function materially changes the
apparent performance and even the ranking of competing models, so metric
selection is itself a substantive modelling decision \parencite{Bruhn2025,STONE1995,SAVABI1995}.
Second, established physically based and analytical models---DRAINMOD, SIDRA-RU,
HYDRUS-2D, and generalisations of the Glover--Dumm and Boussinesq equations---reproduce
observed drainage discharge, water-table profiles, and effluent chemistry with
accuracy that is generally good to excellent once soil hydraulic parameters are
appropriately estimated \parencite{Salazar2008,Jeantet2021,Naderi2025,Mehdinejadiani2013,Mehdinejadiani2020}.
Third, data-driven methods---artificial neural networks (ANNs), random forests
(RF), and deep-learning surrogates---frequently match or outperform conventional
models for the specific outputs they are trained on, while requiring less input
data and computation \parencite{Sarangi2006,Yang2000,Djiemon2019,Ellafi2021,Lei2026}.
Fourth, remote sensing combined with deep learning now maps tile-drained extent
at regional to national scales with accuracies commonly exceeding 85--96\% and
totals within a few percent of agricultural-census estimates \parencite{Valayamkunnath2020,Wan2024,Redoloza2023}.
Fifth, geostatistical (kriging-family) methods deliver reliable, low-bias
estimates of spatially variable soil properties for drainage design and can
reduce sampling effort substantially \parencite{GALLICHAND1991,GALLICHAND1992,Moustafa1998b}.
The following subsections synthesise the quantified outcomes underpinning each
of these lines of evidence.

\subsection{Statistical Metrics and Objective Functions for Model Evaluation}

A central methodological finding is that different performance measures (PMs)
rank models differently, and that this behaviour is systematic rather than
incidental. Across all 14 field datasets analysed by \textcite{Bruhn2025} in
Illinois, the Kling--Gupta efficiency (KGE) consistently outperformed the
Nash--Sutcliffe efficiency (NSE): mismatches between KGE-derived outputs were
always smaller than those between NSE-derived outputs, regardless of the
reference output used, and the same overall PM behaviours recurred across every
dataset, which the authors present as evidence that the conclusions generalise
to field-scale tile-drainage modelling. The heteroscedastic maximum-likelihood
estimator (HMLE) ranked between KGE and NSE overall, but its standing depended
strongly on the underlying error structure: it performed best under scenarios~1
and~2 and worst under scenarios~5 and~6, and in scenario~6 a ranking reversal
occurred for a subset of 4 of 14 datasets (29\%), where HMLE fell to rank~3
while NSE held rank~2 \parencite{Bruhn2025}. In an accompanying DRAINMOD case
study the residuals of neither simulated output showed increasing variance,
leading to the conclusion that KGE would be the best objective function and HMLE
the worst for that situation \parencite{Bruhn2025}. That work also introduced a
quantitative ranking system to summarise PM performance across scenarios and
released open-source HMLE implementations in both R and Python
\parencite{Bruhn2025}.

The limitations of correlation-based statistics are underscored by a paired
discussion of storm-runoff model testing \parencite{STONE1995,SAVABI1995}. The
example simulations grossly overpredicted all 10 events, yet this
unacceptability was not reflected by $R^2$ or the standard error $S_e$; a
regression-based test instead found the intercept significantly different from
zero and the slope significantly different from one, so the model was deemed
unacceptable, and the two null hypotheses were rejected because both
$t$-statistics exceeded 5.041 \parencite{STONE1995,SAVABI1995}. The
corresponding RMSE and mean bias error (24.311~mm and 20.640~mm) were said to
better reflect the high overall differences between simulated and measured
runoff than the correlation measures \parencite{STONE1995,SAVABI1995}.
\textcite{SAVABI1995} note that regression analysis---including tests for
intercept~$=0$ and slope~$=1$ together with the $R^2$ and $S_e$ tests---is a
valid goodness-of-fit approach when the discrepancy between measured and
predicted values is random, and that it can indicate whether model bias is
random or systematic. The interplay of tabulated versus fitted metrics also
appears in curve-number estimation for 12 tile-drained Ohio fields, where
correlation coefficients ranged from 0.29 for the tabular method to 0.55--0.56
for the complacent--violent and $P/S>0.46$ methods, NSE ranged from $-0.26$ for
the maximum-potential-retention method to 0.55 for the complacent--violent
method, and model bias ranged from $-0.05$~mm to $-1.17$~mm; the
complacent--violent method gave the highest correlation and NSE and the best
predictions of the largest outlier events \parencite{Wessel2022}.

\subsection{Performance of Physically Based and Analytical Drainage Models}

\subsubsection{Field-scale water-balance and effluent simulators}

DRAINMOD, evaluated over four years of experiments in south-east Sweden,
achieved good monthly agreement between observed and simulated drainage rates
across five soil-information datasets (H0--H4) \parencite{Salazar2008}. Monthly
mean absolute error (MAE) ranged from 0.57--0.82~cm for conventional drainage
(CD), 0.38--0.41~cm for controlled drainage CWT1, and 0.15--0.22~cm for CWT2;
the modified coefficient of efficiency ($E'$) reached 0.62--0.74 (CD),
0.72--0.74 (CWT1), and 0.79--0.86 (CWT2); and the modified index of agreement
($d'$) reached 0.80--0.93 across treatments \parencite{Salazar2008}. Notably,
substituting ROSETTA-estimated $K_s$ for laboratory-measured values changed the
overall normalised error by less than 3\% in evapotranspiration, less than 1\%
in infiltration, and less than 15\% in drainage; controlled-drainage errors in
predicted drainage were generally under 2\% (except H3 under CWT2 at 15\%),
supporting the conclusion that ROSETTA-estimated $K_s$ can drive DRAINMOD as
accurately as measured values, at least in coarse-textured soils
\parencite{Salazar2008}.

The SIDRA-RU model likewise simulates drainage discharge accurately, especially
on silty soils, with slightly weaker but still acceptable performance on clayey
soils \parencite{Jeantet2021}. For 21 of 22 reference plots the calibration
KGE$'$ exceeded 0.5, spanning 0.54--0.83 on silty soils, 0.54--0.76 on
silty-clay soils, and 0.44--0.76 on clayey soils; annual drained water balances
lay close to the 1:1 line with $R^2>0.9$ for all three textures; low- and
mean-flow biases were near zero while high-flow biases were larger, especially
on clayey soils; and a split-sample test confirmed temporal robustness at the
national scale, enabling long-term impact studies \parencite{Jeantet2021}.
HYDRUS-2D reproduced observed nitrate concentrations in drainage water very well
under both mid-season drainage (MSD) and alternate wetting and drying (AWD) in
Iran, with RMSE of 0.25--0.99~mg~L$^{-1}$, normalised RMSE of 5.3--19.1\%,
normalised mean bias error of $-2.14$ to 10.47\%, and modelling efficiency of
0.6--1.0, with only a slight tendency to overestimate \parencite{Naderi2025}.
The calibrated model estimated the grey water footprint (GWF) far more
consistently with field-based estimates than the tier-1 method: HYDRUS-based GWF
values differed from the field-based method by only 7.0--11.9\% (MSD) and
$-9.4$ to 8.2\% (AWD), whereas the tier-1 method reached relative errors up to
3960.7\%; AWD raised rice yield by about 39.6\% but increased GWF roughly
1042\% relative to MSD \parencite{Naderi2025}.

\subsubsection{Analytical water-table and drain-flow equations}

Generalisations of classical drainage theory have been repeatedly validated
against homogeneous and heterogeneous soils. A space-fractional Boussinesq
formulation reduces to the Glover--Dumm model when the heterogeneity index
approaches 2, so in homogeneous soil the fractional and Glover--Dumm solutions
were nearly coincident with minor error, whereas in a very heterogeneous
experimental field soil (heterogeneity index 1.04) the fractional model tracked
observed water-table profiles more closely and outperformed Glover--Dumm
\parencite{Mehdinejadiani2013}. Extending this to parabolic and elliptical
initial profiles, \textcite{Mehdinejadiani2020} found all fractional models more
sensitive to the heterogeneity parameter $\alpha$ than to $\kappa_x/\mu$;
estimated $\alpha$ was near 2 for relatively homogeneous fine sand but around
1.3 and 1.1 for sandy clay loam and field soils, and the elliptical (SFBE-E) and
first-parabolic--type (SFBE-F) forms generally fit best while SFBE-P2 fit worst,
with SFBE-E substantially improving on its classical counterpart in
heterogeneous soils. A Bees-Algorithm inverse model (DIM) then provided
repeatable, reliable estimates of $\kappa$ (variance across ten runs
$\approx0$), and drainage equations using DIM-derived $\kappa$ gave the best
simulation performance across laboratory and field datasets, outperforming the
Glover--Dumm criterion (GDC), Skaggs's drainage criterion, the drainage-discharge
method, and the direct method; in a worst case at $t=110$~min the global
performance index improved almost 34-fold over the GDC and almost 9-fold over
Skaggs's criterion \parencite{Mehdinejadiani2022}. Complementing these,
\textcite{Samani2007} showed that using water-table-profile variation over time
within an inverse-problem framework let all analytical models agree well with
measured data, that treating effective porosity as a variable improved precision
in every case, that the Topp and Moody model was the most accurate, and that
calibrated $K$ and $f$ values yielded much better water-table predictions than
measured values while remaining close to them. A Green-element method (GEM)
solution of the nonlinear governing equation resolved slope effects efficiently,
matching or bettering earlier analytical results for most practical cases: for
Test~2 GEM gave a maximum error of 2.5\% at 30\% slope, and it predicted better
than Childs and Pavlovsky for a $-15$\% slope, though errors grew as slope
exceeded 0.60 \parencite{Onyejekwe1999}. Field determination of drainage
parameters is also feasible: a Turkish site yielded drainable pore space of
7.0--7.2\%, a mean drainage reaction factor of 0.599~day$^{-1}$, and
transmissivity $K_d$ near 6.7~m$^2$~day$^{-1}$, in agreement with reference
values for medium-textured (silt loam) soils, with the caveat of large
differences among published reference limits \parencite{Bahceci2008}.

\subsubsection{Comparisons of process models with data-driven surrogates}

Head-to-head comparisons show that the best model depends on the target
variable. For an Indian subsurface-drained field, a back-propagation neural
network (BPNN) with 12 hidden neurons and a learning rate of 0.02 predicted
drainage-effluent salinity more accurately than SALTMOD at both 35~m and 55~m
drain spacings, with smaller deviations, but SALTMOD was superior for root-zone
soil salinity because the ANN could not represent the soil-profile salt balance
\parencite{Sarangi2006}. The authors conclude that ANNs are simpler and faster
and operate on minimal input data---moving from a ``black box'' toward a ``gray
box'' when system-based parameters are included---while SALTMOD remains
preferable for detailed salt-balance accounting \parencite{Sarangi2006}.

\subsection{Machine Learning and Data-Driven Predictive Models}

Artificial neural networks have proven accurate and efficient for real-time
water-table management. For Canadian subsurface-drainage (SD) and
subirrigation (SDI) systems, low RMSEs confirmed strong ANN performance: 1996
real-time control gave minimum RMSEs of 73.8~mm (SD) and 60.6~mm (SDI), and
time-series simulation improved for the longer time gap ($D=3$) and with more
training data \parencite{Yang2000}. Most best-performing networks were small,
indicating that large architectures were unnecessary, and the authors emphasise
that ANNs are easy and fast to develop, require little training time, and can be
improved simply by adding recent data \parencite{Yang2000}. Where measured
saturated hydraulic conductivity is scarce, site-specific ANNs predicted $K_{sat}$
more accurately than pedotransfer functions (PTFs) in both an easy-access plain
(EAP, best $R^2=0.90$) and a hard-access plain (HAP, best $R^2=0.96$); PTF-2 was
the best PTF but still less accurate, and ANN-based drainage designs produced
drain spacing and water-table depths close to measured-data designs, whereas
PTF-based designs gave shorter spacing, deeper water tables, and higher initial
costs \parencite{Ellafi2021}. The number of $K_{sat}$ measurements needed to
train an accurate ANN could be reduced to about 50 \parencite{Ellafi2021}.

Random forests outperformed multiple linear regression (MLR) for
drainage-related agronomic prediction in Quebec. For corn plant height, RF
achieved $R^2$ of 0.97 (calibration), 0.84 (validation), and 0.98 (complete
dataset) versus 0.69 for MLR, and RF predicted mean water-table depth with
$R^2=0.95$ versus 0.28 for MLR \parencite{Djiemon2019}. Both methods identified
drain depth and drain distance as important predictors, but RF additionally
captured nonlinear effects and smaller prediction intervals, and it located an
optimal subsurface drain spacing of about 7~m and a critical mean water-table
depth greater than 0.25~m for maximising corn height, while subsoiling direction
had negligible effect \parencite{Djiemon2019}. At the frontier of surrogate
modelling, an Attention U-Net with an adaptive weighted loss (AWL) cut the
average $L_2$ error across four predicted channels by roughly 43\% relative to a
baseline U-Net with mean-squared-error loss, lowering validation loss to
$5.9\times10^{-5}$ versus $8.5\times10^{-5}$ \parencite{Lei2026}. On the test set
the model reached $R^2$ of 0.997--1 and RMSE below 0.05 for hydraulic head,
drainage water volume, salt concentration, and salt discharge; relative errors
of total drainage water and salt discharge stayed within 1.5\%; and a global
optimisation converged to a 40~m drain spacing and 1.4--1.95~m burial depth
while drastically reducing computation time compared with numerical simulation
\parencite{Lei2026}.

Clustering and discriminant analysis provide a further data-driven route to
delineating management zones. In Iowa, average annual NO$_3$-N leaching losses in
drainage water varied from a minimum mean of 6.8~kg-N~ha$^{-1}$ (1994) to a
maximum mean of 35.2~kg-N~ha$^{-1}$ (1993); cluster analysis grouped normalised
losses into high (mean 2.18, 14\% of area), medium (0.28, 50\%), and low
($-0.54$, 36\%) zones, and stepwise discriminant analysis identified elevation
and subsurface drainage as the significant discriminating variables
($p<0.01$), predicting cluster membership with about 86\% accuracy
\parencite{Bakhsh2005b}. A companion yield study formed low, medium, and high
clusters from normalised corn--soybean yields each year and found that soil,
curvature, and aspect contributed significantly in four of six years, though no
variable qualified for the discriminant model in 1997, 1998, or the pooled
1993--1998 data at $P\le0.15$ \parencite{Bakhsh2007b}.

\subsection{Remote Sensing and Deep Learning for Tile-Drainage Mapping}

\subsubsection{Regional and national tile-drainage extent products}

Two large-area products dominate the mapping evidence. AgTile-US, a 30-m
decision-tree map for the contiguous United States, agreed closely with USDA
county estimates, reporting a national tile-drained area of 22{,}340{,}442~ha
against 22{,}418{,}966~ha from USDA (a $-0.35$\% difference); validation with
16{,}000 ground-truth points gave 86.03\% overall accuracy (13.97\% omission
error), with state-level accuracy highest in the heavily drained north-central
states (72.8--100\%) \parencite{Valayamkunnath2020}. The complementary
random-forest product SEETileDrain, built through Google Earth Engine for the US
Midwest, achieved testing overall accuracy of 0.96, recall 0.96, precision 0.85,
balanced accuracy 0.96, and F1 of 0.90, with county-scale area agreeing with the
Ag~Census at $r^2=0.69$; it estimated 185{,}549~km$^2$ of tile drainage in 2017,
about 9.4\% below the USDA-NASS estimate of 204{,}842~km$^2$
\parencite{Wan2024}. A simplified random forest using only slope and soil
drainage classes performed far worse (out-of-bag error 36\%, precision 0.29, F1
0.4), and 27\% of counties with NASS-reported tile were predicted to have none,
though these had a small median reported area of about 27~km$^2$
\parencite{Wan2024}. In Denmark, a selective-ensemble machine-learning approach
predicted artificially drained areas over 49.4\% of the study region, with the
final ensemble achieving $\kappa=0.526$ and 76.5\% simple accuracy; models
trained on principal components outperformed those on raw covariates and had
shorter prediction times (mean 8.2~s in the ensemble versus 49.4~s in the
library), and evaluation performance correlated with hillclimb performance
($R^2=0.66$) \parencite{Moller2018}.

\subsubsection{Deep-learning detection of drain lines from imagery}

Convolutional segmentation networks reliably delineate visible drain lines. A
U-Net trained on imagery from the US Midwest delineated tile drains in 53
validation images across two areas---including 21 images about 200~km southwest
of the training data---exceeding 92\% accuracy and roughly 53\% balanced
accuracy (reported 93--96\% accuracy, 52--54\% balanced accuracy), with training
performance of 95\% accuracy and 63\% balanced accuracy; the model learned to
ignore roads, buildings, trees, lakes, and boundary artefacts, and drains were
most visible in spring imagery under high solar radiation, intermediate soil
water, and minimal crop cover \parencite{Redoloza2023}. In Ohio, a
deep-learning segmentation model reached minimum Dice losses of 0.28 (training)
and 0.42 (validation) over 500 epochs, an average validation Dice coefficient of
0.58 across 257 images, and 0.55 on a large-scale map, outperforming Sobel and
Canny edge detectors that produced noisier and more false detections
\parencite{Song2021}. An ensemble fully convolutional network gave accurate,
robust drainage-pipe detection on training, validation, and unseen test data,
was visually closer to the ground truth than individual models, and outperformed
Sobel edge detection and the Hough transform; performance depended on input
resolution, peaking at 400~m/256~pixels with reasonable results (Dice $>0.3$)
between 200 and 500~m/256~pixels \parencite{Woo2023}. Earlier automated
approaches were already useful: a decision-tree classification of an Indiana
watershed identified 79\% as potentially tile-drained and, after masking, an
estimated 50\% drained area with drain spacing of 17--80~m and sub-watershed
averages of 43--44~m, within the recommended 23--45~m range for local soils
\parencite{Naz2009}.

\subsubsection{Ancillary remote-sensing methods relevant to drainage}

Several studies quantify how remote sensing infers surface conditions tied to
drainage. Synthetic-aperture-radar (SAR) image simulation of undulated corn
canopy reproduced the statistical and undulation properties of a validation
parcel, with autocorrelation peaks separated by 5.3~pixels (25~m) corresponding
to distinct drainage ridges; coefficient of variation and contrast discriminated
corn fields with and without subsurface drainage, and low contrast (near-zero
plant-height differential) was associated with drained fields
\parencite{Meunier2001}. A Landsat-based soil-water-content algorithm exploited
strong SWIR absorption (Bands~5 and~7), with reflectance nearly constant below
0.20~cm$^3$~cm$^{-3}$ and rising above 0.40~cm$^3$~cm$^{-3}$; the best 15-cm
model, $Y_{15}=0.48e^{-2.64(B5-B1)}$, validated on 70 observations across nine
fields with $r=0.94$ and RMSE near 0.02~cm$^3$~cm$^{-3}$, though it
underestimated wetter conditions \parencite{Rijal2013}. For cover-crop
monitoring in an Illinois watershed, SAVI, EVI, and TVI correlated strongly with
biomass and N uptake; SAVI was selected because it matched or bettered the
others while needing only red and near-infrared bands, giving biomass estimates
of 86--107\% and N-uptake estimates of 85--98\% of ground values, and it
outperformed ordinary kriging and inverse-distance weighting where management
varied \parencite{Xu2018}. A large-aperture-scintillometer study showed that
sensible-heat-flux retrieval is sensitive to parameterisation: a rule-of-thumb
roughness length overestimated $H_{LAS}$ by 47--66\% via inflated $u^*$, but an
improved $z_0$ estimate brought $u^*$ RMSD from 0.21 to 0.08~m~s$^{-1}$ and slope
from 1.37 to 0.98, yielding LAS--eddy-covariance agreement within 32~W~m$^{-2}$
and a slope of 1.0 for the best (W71) temperature function \parencite{Zhang2010}.
A broader review documents an increasing 20-year trend in agricultural remote
sensing, a sharp post-2015 rise in unmanned aerial systems (34\% of papers from
Europe, 20\% USA, 11\% China), and notes that high-resolution thermal imagery
can detect temperature differences over drain lines and thereby the subsurface
tile network, while subsurface drainage remained comparatively under-studied
\parencite{Khanal2020}.

\subsection{Geostatistical and Spatial-Interpolation Methods for Drainage Design}

Kriging-family methods are consistently the most reliable interpolators for the
spatially variable soil properties that govern drain spacing. In the Nile Delta,
variations in hydraulic conductivity were spatially isotropic with a high nugget
and a 5.6~km range of influence; kriging cross-validation gave low estimation
errors (mean $-4.2\times10^{-4}$, variance 1.012), and kriged effective
conductivities were recommended as design values because they are the best
obtainable estimates and produce smooth surfaces delineating homogeneous blocks
\parencite{GALLICHAND1991}. Grid-optimisation work showed that a 600~m grid
retained nearly all information: average kriged point conductivity was
0.201~m~d$^{-1}$ on both complete and reduced grids, block-effective correlations
reached 0.906--0.909 with average absolute deviations of 0.028--0.031~m~d$^{-1}$
(versus 0.930 and 0.031 for the measured grid), and a preliminary 900~m survey
sufficed to fix the optimal sampling density \parencite{GALLICHAND1992}. For clay
content, no interpolation method showed important bias (estimated averages
56.19--57.07\% versus 57.13\% observed); cokriging gave the lowest mean absolute
difference (12.01) but only marginally beat kriging and inverse-distance
weighting, while the closest-neighbour method gave the highest MAD (16.44) and
the largest low-clay areas, so kriging, cokriging, and weighted moving average
($\alpha=1$) were all judged adequate \parencite{GALLICHAND1993}.

The practical value of accounting for spatial variability is quantified by
several Egyptian studies. Kriged estimates outperformed measured point values:
measured water-table depth underestimated the need for subsurface drainage by
about 40\%, measured soil salinity overestimated levels by 59\%, and using
kriged hydraulic conductivity widened calculated lateral drain spacing by about
5\% (up to 9.6\% for some unsteady-state formulas), saving 17{,}155~m of lateral
pipe while reducing the minimum sample size roughly three- to sevenfold
\parencite{Moustafa1998b}. A geostatistical correlation model for representative
$K_{sat}$ was validated on Nile Delta and Nile Valley soils and shown practically
valuable; neglecting spatial variability of $K_{sat}$ could bias design drain
spacing by $-27$\% to $+3$\% and overestimate required sample size by about 76\%
\parencite{Moustafa2000}. The covariance variogram proved superior to the
semi-variogram for clay soils, being smoother, more robust to high values, and
lower in cross-validation error, and subsurface drainage itself strongly altered
spatial variability---the spatial correlation range increased about 29\% and
relative structured heterogeneity doubled four years after installation, while
the nugget rose for all soil-water properties \parencite{Moustafa1998}. In
Turkey, ordinary point kriging was appropriate for $K_{sat}$ and soluble salts,
which showed strong depth-dependent structure (horizontal models explaining
81--84\% of variance) with horizontal correlation distances of 710--1130~m,
supporting an approximately 1130~m~$\times$~1130~m subfield and $\geq75$~cm profile
as a reference for drain-tile layout \parencite{Akis2015}. Simple regression and
geostatistical models for rapid representative $K_{sat}$ were also proposed for
large-scale drainage design, with coefficients that necessarily differ among
soils and hydrological conditions \parencite{Sepaskhah2004}.

\subsection{Calibration, Parameter Estimation, and Inverse Modelling}

The way objective functions are constructed and optimised strongly affects
parameter identifiability. In Iranian inverse modelling of soil hydraulic
parameters, a genetic algorithm (GA) generally outperformed PEST, especially for
objective functions using two or three datasets; combining drainage discharge,
water-table depth, and drainage salinity was the best strategy, though drainage
discharge with water-table depth alone already gave excellent results, and PEST
was less capable for complex objective functions \parencite{Sedaghatdoost2019}.
Variational data assimilation (VDA)---the first such implementation for
SIDRA-type models---used global adjoint sensitivity analysis to identify
$K_{sat}$ as the most influential parameter (followed by $\mu$ then $S_{inter}$),
selecting only 4 of 11 parameters for calibration; split-sample validation gave
KGE $>0.66$, NSE $>0.53$, and RMSE $<2.44$~l~s$^{-1}$, and VDA cut the five-year
cumulative-volume discrepancy to 4.8~mm ($<1$\%) versus 58.7~mm (10\%) for
PAP-GR \parencite{Chelil2022b}. Single-parameter tests showed $K_{sat}$ and
$\mu$ converged uniquely while $SSDI$ and $S_{inter}$ hit local minima, and a
hybrid background-variance approach improved robustness across periods
\parencite{Chelil2022b}. At regional scale, a single-objective calibration in
California produced identifiable parameters with small uncertainties and
predicted regional water-table variation well (RMSE 2.76~m, correlation 0.98),
while multi-objective calibration revealed clear trade-offs between water-table
and pumping predictions; drainage predictions were much more sensitive yet
largely independent of the other objectives, and prediction uncertainty was
large for water tables but small for drainage and pumping \parencite{Schoups2005}.

\subsection{Evaluation of Measurement, Sampling, and Sensor Methods}

Performance evaluation extends to the field protocols that generate the data
against which models are judged. For phosphorus-load monitoring in Michigan tile
drainage, all sampling strategies underestimated total reactive phosphorus (TRP)
load, but the magnitude depended sharply on design: time-proportional discrete
sampling worsened from $-0.2$\% at 3~h to $-51.0$\% at 14~d, and time-proportional
composites ran from $-12.4$\% (1-d) to $-43.0$\% (7-d), whereas flow-proportional
discrete sampling stayed within $-0.2$ to $-5.4$\% and flow-proportional
composites within $-0.5$ to $-7.2$\% \parencite{Dialameh2022}. Flow-proportional
strategies thus gave more accurate cumulative-load estimates with fewer samples
and lower analysis cost, and the number of aliquots did not substantially affect
accuracy \parencite{Dialameh2022}. A recording ultrasonic flow meter in a
multi-slotted U-pipe achieved very high calibration quality (mean and minimum
$R^2$ of 0.9993 and 0.993), propagated $Q(h)$ uncertainties of only 0.1--0.4\%
across slot sizes, and performed well in laboratory and field with good signal
stability using low-cost components \parencite{Enamorado2007}. For denitrifying
bioreactor tracer tests, a continuously deployed bromide sensor produced both
under- and over-reported concentrations---the best field match (test~\#7) still
underpredicted the peak by 7.5\%, with lowest error RMSE 5.0~mg~L$^{-1}$ and MAE
3.0~mg~L$^{-1}$ and a pooled slope near 1.0 (0.96); most tracer metrics fell
within $\pm10$\% of lab-based values, but recovery and peak-timing metrics could
change substantially (e.g.\ hydraulic efficiency from 0.29 to 1.8), and post-hoc
bottle-based measurement outperformed real-time deployment (pooled slope 0.97,
RMSE 3.1~mg~L$^{-1}$) \parencite{deOliveira2023b}. A non-intrusive geophysical
protocol produced ``as-built'' subsurface maps of golf-green tile, confirming
all tile spacing met the USGA specification but finding four of 12 laterals below
the minimum required downward slope (two sloping upward) and about 20\% of the
area failing root-zone thickness standards \parencite{Freeland2014}. On
drainage-efficiency measurement itself, both large-ring and small-core methods
were valuable for relative comparisons, the large ring covering a more
representative elementary volume, and infiltration and saturated conductivity
were 5--65 times larger in the trench backfill zone than midway between trenches
in Swedish clays \parencite{Messing2006}.

\subsection{Environmental Indicators and Empirical Regression Relationships}

Aggregative indicators and empirical regressions provide a complementary,
policy-facing layer of model evaluation. A global compilation of 60
hydraulic-retention-time (HRT)--nitrate-removal-efficiency (NRE) regressions
showed strong within-study linearity but wide between-study variability (slopes
0.13--18~\%/h, intercepts $-19$ to $+47$\%); slope correlated positively with
water temperature (Spearman $\rho=+0.39$) and negatively with inflow N
concentration ($\rho=-0.46$), while the intercept correlated with no compiled
parameter \parencite{Christianson2025c}. Generalised median models
were NRE~$=2.9\times\mathrm{HRT}+0.37$ for the full dataset and
$3.2\times\mathrm{HRT}+0.18$ for a subsurface-drainage ``Tile Proxy'' subset, and
study scale did not significantly alter slope or intercept, though field-scale
studies showed greater slope heterogeneity \parencite{Christianson2025c}. A soil
quality index (SQI) built by scoring-function analysis in Ohio ranged narrowly
(0.69--0.71) and was not significantly affected by tillage or subsurface drainage
but was significantly correlated with corn yield ($R=0.62$), with $K_{sat}$ the
key physical indicator and soil organic carbon the highest-weighted chemical
indicator \parencite{Nakajima2015}. At landscape scale, phosphorus-loss
indicators such as IROWC-P are framed as aggregations of multiple source $\times$
transport interactions whose dominant components depend on climate, soils,
topography, and management; the revised component structure estimates
edge-of-field loss and applies a delivery ratio, and results are expected to be
``directionally and magnitudinally correct'' against annual water-quality data
\parencite{Reid2018}. More broadly, environmental indicators are presented as
powerful tools for tracking change and informing policy that function well at
field, catchment, and local scales but need refinement for national or
international comparison, with the Phosphorus Index highlighted as a widely
adopted aggregative indicator encompassing subsurface tile-drainage pathways and
soil organic carbon as an indirect indicator of soil quality and CO$_2$ emissions
\parencite{Zhang2018}. Finally, a riparian-quality assessment of 21 tile-drained
Canadian sites produced RQI scores of 45--129 that correlated with an independent
index ($R^2=0.56$), while width alone correlated more strongly with that index
($R^2=0.79$), and partial-least-squares analysis showed nutrient and sediment
retention related differently to landscape variables by region
\parencite{Terrado2015}.

\printbibliography
\end{document}
